% =====================================================================
%  第 6 章  当前进度与后续计划
% =====================================================================

\section{当前进度与后续计划}
\label{sec:progress}

\subsection{已完成的工作}

目前完成的工作包括以下五项。

\begin{enumerate}
  \item \textbf{选题论证}。确定了研究对象（极值图论中的不变量不等式）、
  核心判断（主要困难位于筛选环节）以及三个子问题的划分。
  \item \textbf{数据构造}。完成 27810 张图的数据集，九个不变量全部精确计算，
  并以四条已知结论作校验，违反数为零（见表\ref{tab:validate}）。
  \item \textbf{核心算法}。完成候选形式的表示与线性规划求解、三类反例搜索、
  置换检验与 BH 校正。线性规划已能复现五条已知结论。
  \item \textbf{基线实验}。在不调用模型的条件下跑通生成与筛选全流程，
  并完成反例搜索策略的对照实验（第\ref{sec:results}节）。
  \item \textbf{系统实现}。智能体主循环、模型客户端、提示词与轨迹记录均已实现，
  演练模式通过。密钥管理与版本控制亦已配置。
\end{enumerate}

另有一项早期工作值得说明。在确定选题之前，实现了一个极简的 ReAct 循环
（思考、行动、观察），用于理解智能体的工作机制。
该原型仅含两个示例工具，与本课题的数学内容无关，
但验证了“模型能够依据工具返回的结果决定下一步动作”这一基本机制。

\subsection{待完成的工作与时间安排}

后续工作分为四个阶段，如表\ref{tab:stages} 所示。

\begin{table}[H]
\centering
\small
\caption{后续工作阶段划分}
\label{tab:stages}
\begin{tabular}{@{}clp{7.4cm}@{}}
\toprule
阶段 & 内容 & 产出 \\
\midrule
一 & 接入模型，跑通完整闭环 & 至少三轮的轨迹日志与反例解释 \\
二 & 消融实验 & 模型提议与随机提议的对照数据 \\
三 & 筛选的对照实验 & 校正前后的无效命题比例 \\
四 & 分析与撰写 & 报告后半部分的定稿 \\
\bottomrule
\end{tabular}
\end{table}

各阶段的日程安排如表\ref{tab:schedule} 所示。

\begin{table}[H]
\centering
\small
\caption{时间安排}
\label{tab:schedule}
\begin{tabular}{@{}llp{5.4cm}l@{}}
\toprule
日期 & 内容 & 产出 & 状态 \\
\midrule
11/19 & 选题论证、报告框架 & 报告前两章 & 完成 \\
11/20 & 数据、算法与基线实验 & 数据集、图表与第\ref{sec:results}节 & 完成 \\
11/21 & 接入模型，跑通闭环 & 轨迹日志 & 待办 \\
11/22 & 消融实验与筛选对照 & 对照数据 & 待办 \\
11/23 & 报告补充与图表整理 & 报告定稿 & 待办 \\
11/24 & 汇报材料 & 幻灯片 & 待办 \\
11/25 & 中期汇报 & --- & --- \\
\bottomrule
\end{tabular}
\end{table}

\subsection{已知风险}

已知风险及相应对策如表\ref{tab:risks} 所示。

\begin{table}[H]
\centering
\footnotesize
\caption{已知风险与对策}
\label{tab:risks}
\begin{tabular}{@{}p{3.9cm}p{10.1cm}@{}}
\toprule
风险 & 对策 \\
\midrule
不变量计算过慢 & 色数与独立数均为 NP-难问题，大规模图上无法精确计算。
对策：小规模图精确计算，中等规模图以上下界代替并明确标注，
大规模图仅作采样。宁可结论保守，不作精确性上的虚假陈述。 \\
\addlinespace
大语言模型调用不稳定 & 全部结果落盘缓存，重跑时读取缓存；
汇报以录屏备份，不作现场演示。 \\
\addlinespace
模型贡献有限 & 第\ref{sec:method}节已设计消融实验。
若模型与随机选择的产出无显著差异，则如实报告这一结论。 \\
\addlinespace
时间不足 & 优先保证闭环可运行、轨迹可展示；
筛选的对照实验可缩减规模。 \\
\addlinespace
结果不理想 & 失败模式本身即为结论，
如实报告较之虚构成功更具价值。 \\
\bottomrule
\end{tabular}
\end{table}
